% problem-type: truyen-song-3d-mien-anh-huong-t0
% order: 42
% Nguon: De thi MAT3365 (PTDHR Toan tin), cuoi ky, K68, Cau 3. (co dap an chinh thuc)

\section{Truyền sóng 3D có nguồn --- miền ảnh hưởng \& thời điểm rời cân bằng $t_0$}

\subsection*{Đề bài ví dụ}

Xét bài toán Cauchy cho phương trình truyền sóng
\[
  u_{tt}(x,y,z,t)=4\,\Delta u(x,y,z,t)+f(x,y,z,t),\quad (x,y,z)\in\mathbb{R}^3,\ t>0,
\]
với điều kiện ban đầu $u(x,y,z,0)=0$ và $u_t(x,y,z,0)=0$, trong đó
\[
  f(x,y,z,t)=\chi_{[2024,2026]}(x)-\chi_{[0,1]}(t)\,\chi_B(x,y,z),
\]
$B=\{(x,y,z):(x-2025)^2+y^2+z^2\le1\}$.

\begin{enumerate}
\item[(a)] Tìm thời điểm sớm nhất $t_0$ mà $u(0,0,0,t)$ chuyển từ trạng thái cân bằng sang không cân bằng.
\item[(b)] Tính $u(0,0,0,t)$ theo $t$.
\end{enumerate}

\emph{\small(Nguồn: Đề thi MAT3365 --- PTĐHR Toán tin, cuối kỳ, K68, Câu 3. Có đáp án chính thức.)}

\subsection*{Nhận ra dạng này khi nào?}

\begin{itemize}
\item Sóng 3D có nguồn, hỏi giá trị tại \textbf{một điểm} ở \emph{xa} vùng nguồn $+$ hỏi \textbf{thời điểm bắt đầu} có tín hiệu.
\item Nguồn \textbf{trộn}: một phần phụ thuộc một biến phẳng ($x$), một phần đối xứng cầu quanh một tâm khác gốc.
\end{itemize}

\begin{luuybox}
\textbf{``Trạng thái cân bằng''} $=$ trạng thái nghỉ $u\equiv0$. Dữ liệu ban đầu $u=u_t=0$ nên lúc đầu mọi điểm ở cân bằng.
Vì sóng \textbf{lan với tốc độ hữu hạn} $c$ (ở đây $c^2=4\Rightarrow c=2$), một điểm $P$ vẫn $u(P,t)=0$ cho tới khi tín hiệu từ vùng nguồn kịp tới. Do đó
\[
  t_0=\frac{\operatorname{dist}\big(P,\ \operatorname{supp} f\big)}{c}.
\]
\end{luuybox}

\subsection*{Công thức \& phương pháp}

\subsubsection*{Chồng chất: tách nguồn}
% method: nguyen-ly-chong-chat
\begin{dieukienbox}
Dùng khi: PDE \textbf{tuyến tính} --- tách bài toán thành nhiều phần đơn giản hơn,
giải từng phần rồi cộng lại.
\end{dieukienbox}

\begin{nhobox}
\textbf{Nguyên lý chồng chất}: với $u_{tt}-\Delta u = f$, $u(0)=g$, $u_t(0)=h$, tách theo VAI TRÒ:
\[
  u = \underbrace{u^{(g)} + u^{(h)}}_{\text{phần do DỮ LIỆU (nguồn tắt)}} + \underbrace{u^{(f)}}_{\text{phần do NGUỒN (dữ liệu tắt)}}.
\]
\begin{itemize}
  \item $u^{(g)}$: $u_{tt}=\Delta u$ (không nguồn), $u(0)=g$, $u_t(0)=0$.
  \item $u^{(h)}$: $u_{tt}=\Delta u$ (không nguồn), $u(0)=0$, $u_t(0)=h$.
  \item $u^{(f)}$: $u_{tt}-\Delta u=f$ (có nguồn), $u(0)=0$, $u_t(0)=0$ --- dùng Duhamel.
\end{itemize}
\textbf{Luôn kiểm tra}: cộng các bài con phải ra lại đúng PDE $+$ dữ liệu của bài gốc.
\end{nhobox}

\begin{meobox}
Nếu $h = h_1 - h_2$ (hiệu hai hàm), tách tiếp: $u^{(h)} = u^{(h_1)} - u^{(h_2)}$.
Ví dụ: $h = \chi_{B_1(0,0,2)} - \chi_{B_1(0,0,-2)}$ $\Rightarrow$ hai bài toán con, xử lý từng quả cầu.
\end{meobox}


\subsubsection*{Nguyên lý Duhamel (phần nguồn)}
% method: nguyen-ly-duhamel
\begin{dieukienbox}
Dùng khi: phương trình sóng \textbf{có nguồn} $u_{tt}-c^2\Delta u=f$,
dữ liệu ban đầu bằng $0$ (hoặc đã tách phần do dữ liệu ban đầu bằng chồng chất).
\end{dieukienbox}

\begin{nhobox}
\textbf{Định nghĩa $w(x,t;s)$:} với mỗi $s$ cố định, $w$ giải phương trình \emph{không có nguồn}
$w_{tt}=c^2\Delta w$, nhưng đặt dữ liệu tại $t=s$ với \textbf{vận tốc đầu $=f(\cdot,s)$} (lát cắt nguồn):
\[
  \begin{cases}
    w_{tt} = c^2\Delta w & (t>s),\\
    w(x,s;\,s) = 0 & (\text{vị trí đầu}=0),\\
    w_t(x,s;\,s) = f(x,s) & (\text{vận tốc đầu}=f\neq0).
  \end{cases}
\]
\textbf{Nguyên lý Duhamel} --- gộp theo $s$:
\[
  u(x,t) = \int_0^t w(x,t;\,s)\,ds .
\]
\end{nhobox}

\begin{nhobox}
\textbf{Tìm $w$ thế nào:} dùng đúng công thức nghiệm \emph{không nguồn} đã biết (Kirchhoff/d'Alembert),
với dữ liệu $g=0$, $h=f(\cdot,s)$, và thay thời gian $t\to t-s$ (vì $w$ chỉ phụ thuộc thời gian trôi $t-s$).
Với vận tốc lan truyền $c$ (phương trình $u_{tt}=c^2\Delta u+f$):
\[
  \text{1D (d'Alembert): }\ w=\frac{1}{2c}\!\int_{x-c(t-s)}^{x+c(t-s)}\! f(\xi,s)\,d\xi,
\]
\[
  \text{3D (Kirchhoff): }\ w=\frac{1}{4\pi c^2(t-s)}\int_{\partial B_{c(t-s)}(x)}\! f(y,s)\,dS(y).
\]
(Hệ số vận tốc là $\tfrac{1}{2c}$, \emph{không} phải $\tfrac12$; khi $c=1$ thì hai cái trùng nhau.)
\end{nhobox}

\begin{luuybox}
\textbf{Đừng nhầm hai chữ ``không nguồn''!}
\begin{itemize}
  \item \emph{Phần do dữ liệu} $u^{(g)},u^{(h)}$ (chồng chất): không nguồn $+$ dữ liệu THẬT $(g,h)$. Có thể $=0$ ở vài điểm (vd do đối xứng) --- không liên quan Duhamel.
  \item \emph{Sóng con $w$} (Duhamel): không nguồn nhưng dữ liệu đầu $=f(\cdot,s)\neq0$, nên $w\neq0$. Đây mới là thứ Duhamel cộng lại.
\end{itemize}
``Không nguồn'' nói về \textbf{phương trình} (vế phải $=0$), KHÔNG phải dữ liệu $=0$.
\end{luuybox}

\begin{meobox}
$f$ có giá compact theo $s$ (vd $s\in[0,1]$) $\Rightarrow$ tích phân Duhamel tự thu hẹp về đoạn đó.
\end{meobox}


\subsubsection*{Hạ thấp chiều \& đổi biến cầu}
% method: phuong-phap-ha-thap
\begin{dieukienbox}
Dùng khi: dữ liệu hoặc nguồn \textbf{không phụ thuộc một hoặc vài biến không gian}.
Ví dụ: nguồn $f(x,t)=\chi_{[-1,1]}(x_1)$ không phụ thuộc $y,z$
$\Rightarrow$ bài toán 3D thực ra chỉ ``sống'' trong 1D theo $x_1$.
\end{dieukienbox}

\begin{nhobox}
\textbf{Ý tưởng hạ thấp}: nếu mọi thứ chỉ phụ thuộc $x_1$ (không có $y,z$),
thì $\Delta u = u_{x_1 x_1}$ và bài toán trở thành phương trình sóng \textbf{1D}.

Áp dụng công thức d'Alembert cho bài 1D, hoặc suy nghiệm 2D/1D từ công thức 3D (Kirchhoff) bằng cách tích phân thêm theo biến phụ.
\end{nhobox}

\begin{meobox}
Dấu hiệu nhận ra: $f$ hoặc $h$ viết theo $\chi_{[-1,1]}(x_1)$ --- chỉ có $x_1$, không có $y,z$.
Khi đó có thể bỏ hai biến $y,z$ và giải bài 1D nhanh hơn nhiều.
\end{meobox}

\emph{(Nguồn chỉ phụ thuộc $x$ $\to$ sóng 1D + d'Alembert. Nguồn đối xứng cầu quanh tâm $O'$ $\to$ đặt $v=\rho\,U$ với $\rho=$ khoảng cách tới $O'$; giá trị tại điểm cách $O'$ một đoạn $R$ là $U(R,t)=v(R,t)/R$.)}

\subsubsection*{Miền phụ thuộc / vận tốc hữu hạn}
% method: phan-tich-mien-phu-thuoc
\begin{dieukienbox}
Dùng khi: cần xác định nghiệm \emph{có bằng $0$ hay không} tại một điểm xa nguồn/dữ liệu.
Áp dụng cho mọi PDE hyperbolic vận tốc hữu hạn ($c=1$ ở đây).
\end{dieukienbox}

\begin{nhobox}
\textbf{Quy trình 3 bước:}
\begin{enumerate}
  \item Xác định giá (\emph{support}) của $g$, $h$, $f$.
  \item Tính khoảng cách từ điểm quan sát $x_0$ đến giá:
        $d_{\min} = \operatorname{dist}(x_0, \operatorname{supp})$.
  \item Kết luận: $u(x_0,t)=0$ khi $t < d_{\min}$ (sóng chưa tới);
        tín hiệu xuất hiện từ $t = d_{\min}$.
\end{enumerate}
Trong 3D (nguyên lý Huygens mạnh): tín hiệu cũng tắt sau $t = d_{\max}$
($d_{\max}$: khoảng cách lớn nhất đến giá).
\end{nhobox}

\begin{meobox}
Ví dụ: $x_0=(100,0,0)$, $\operatorname{supp}(h)\subset B_1(0,0,\pm 2)$.
Khoảng cách gần nhất $\approx 97$, xa nhất $\approx 103$.
$\Rightarrow$ $u^{(h)}(x_0,t)=0$ ngoài khoảng $t\in[97,103]$ (tính chính xác theo hình học).

Cho nguồn $f=\chi_{[0,1]}(s)\chi_{[-1,1]}(x_1)$: $\operatorname{dist}(x_0,\operatorname{supp}(f(\cdot,s)))\ge 99$,
cần $t-s\ge 99$ $\Rightarrow$ tín hiệu từ $f$ bắt đầu từ $t\approx 99$.
\end{meobox}


\subsection*{Đề bài \& bài giải}

\paragraph{(a) Thời điểm rời cân bằng $t_0$.}\leavevmode

\emph{Bước a.1 --- Nhắc lại khái niệm.}
\begin{itemize}
\item \textbf{Trạng thái cân bằng} của bài này $=$ trạng thái nghỉ $u\equiv0$. Vì dữ liệu ban đầu $u(\cdot,0)=0$ và $u_t(\cdot,0)=0$ nên tại $t=0$ mọi điểm đều cân bằng.
\item \textbf{Vận tốc truyền sóng hữu hạn:} phương trình $u_{tt}=c^2\Delta u+f$ truyền tín hiệu với vận tốc $c$. Ở đây $c^2=4\Rightarrow c=2$. Trước khi tín hiệu kịp tới, điểm hỏi vẫn ở trạng thái cân bằng.
\item \textbf{Miền ảnh hưởng (forward cone)} của một điểm nguồn $Q$ phát tại $s\ge0$ là $\{(P,t):|P-Q|\le c(t-s)\}$. Hợp các nón này theo $(Q,s)\in\operatorname{supp} f$ là tập các điểm $(P,t)$ \emph{có thể} bị nguồn tác động.
\item Hệ quả: $u(P,t)=0$ khi nào nón ngược $\{(Q,s):|Q-P|\le c(t-s),\ 0\le s\le t\}$ chưa giao $\operatorname{supp} f$, tức $c\,t<\operatorname{dist}(P,\operatorname{supp} f)$.
\end{itemize}

\emph{Bước a.2 --- Tính $\operatorname{dist}(P,\operatorname{supp} f)$.} Điểm hỏi $P=(0,0,0)$. Giá $\operatorname{supp} f$ gồm hai phần:
\begin{itemize}
\item Dải phẳng $S_1=\{(x,y,z):2024\le x\le2026\}$: điểm gần $P$ nhất nằm trên mặt $x=2024$, khoảng cách $=2024$.
\item Quả cầu $B$ tâm $O'=(2025,0,0)$ bán kính $1$: điểm gần $P$ nhất nằm trên đoạn $\overline{PO'}$, cách $P$ một khoảng $|PO'|-1=2025-1=2024$.
\end{itemize}
Vậy $\operatorname{dist}(P,\operatorname{supp} f)=\min(2024,2024)=2024$.

\emph{Bước a.3 --- Suy ra $t_0$.}
\[
  t_0=\frac{\operatorname{dist}(P,\operatorname{supp} f)}{c}=\frac{2024}{2}=1012.
\]
Với $0\le t<1012$ thì $u(0,0,0,t)=0$ (cân bằng); tại $t_0=1012$ điểm $P$ rời cân bằng.

\paragraph{(b) Bước 0 --- chồng chất.}\leavevmode

\emph{Bước 0.1 --- Tách nguồn.} Viết $f=f_1-f_2$ với
\[
  f_1(x,y,z,t)=\chi_{[2024,2026]}(x),\qquad
  f_2(x,y,z,t)=\chi_{[0,1]}(t)\,\chi_B(x,y,z)=\chi_{[0,1]}(t)\,\chi_{[0,1]}(\rho),
\]
trong đó $\rho:=\sqrt{(x-2025)^2+y^2+z^2}$ là khoảng cách tới $O'=(2025,0,0)$. Phần $f_1$ chỉ phụ thuộc $x$ (\emph{đối xứng phẳng}); phần $f_2$ \emph{đối xứng cầu} quanh $O'$.

\emph{Bước 0.2 --- Vì sao $u=u_1-u_2$?} Xét hai bài toán phụ
\[
  (1)\ \begin{cases}u_{1,tt}=4\Delta u_1+f_1\\ u_1(\cdot,0)=u_{1,t}(\cdot,0)=0\end{cases}
  \qquad
  (2)\ \begin{cases}u_{2,tt}=4\Delta u_2+f_2\\ u_2(\cdot,0)=u_{2,t}(\cdot,0)=0\end{cases}
\]
Vì toán tử $\partial_t^2-4\Delta$ \textbf{tuyến tính}, đặt $\tilde u:=u_1-u_2$ thì
\[
  \tilde u_{tt}-4\Delta\tilde u=(u_{1,tt}-4\Delta u_1)-(u_{2,tt}-4\Delta u_2)=f_1-f_2=f,
\]
và $\tilde u(\cdot,0)=0$, $\tilde u_t(\cdot,0)=0$. Bài Cauchy đang xét có nghiệm \textbf{duy nhất} (sóng tuyến tính trên $\mathbb{R}^3$), nên $u=\tilde u=u_1-u_2$. Việc tách như vậy có lợi vì mỗi $u_j$ hưởng \emph{một loại đối xứng} dễ khai thác.

\paragraph{(b) Bước 1 --- $u_1$ (nguồn phẳng $\to$ 1D).}\leavevmode

\emph{Bước 1.1 --- Hạ thấp chiều.} Nguồn $f_1=\chi_{[2024,2026]}(x)$ chỉ phụ thuộc $x$ và dữ liệu đầu $=0$. Do tính \textbf{duy nhất} nghiệm + bất biến phương trình theo phép tịnh tiến $(y,z)$, nghiệm $u_1$ cũng không phụ thuộc $y,z$. Đặt $u_1(x,y,z,t)=v_1(x,t)$ thì $\Delta u_1=v_{1,xx}$, nên
\[
  v_{1,tt}=4\,v_{1,xx}+\chi_{[2024,2026]}(x),\qquad v_1(x,0)=v_{1,t}(x,0)=0,\quad x\in\mathbb{R}.
\]

\emph{Bước 1.2 --- Duhamel 1D. ``$s$'' là gì?} \textbf{Nguyên lý Duhamel} biến bài có nguồn thành tích phân theo \emph{thời điểm phát} $s\in(0,t)$ của \emph{sóng con thuần nhất} $w(x,t;s)$ — nghiệm bài $w_{tt}=4w_{xx}$ với dữ liệu tại $t=s$ là $w(\cdot,s)=0$, $w_t(\cdot,s)=f_1$. Công thức d'Alembert 1D với $c=2$ cho
\[
  w(x,t;s)=\frac{1}{2c}\int_{x-c(t-s)}^{x+c(t-s)}\!\chi_{[2024,2026]}(\xi)\,d\xi,
\]
(miền phụ thuộc là khoảng độ dài $2c(t-s)$ vì sóng con đã lan trong thời gian còn lại $t-s$). Cộng lại:
\[
  v_1(x,t)=\int_0^t w(x,t;s)\,ds
  =\frac{1}{2c}\iint_{\Delta(x,t)}\!\chi_{[2024,2026]}(\xi)\,d\xi\,ds,
\]
với $\Delta(x,t)=\{(s,\xi):0<s<t,\ x-c(t-s)<\xi<x+c(t-s)\}$.

\emph{Bước 1.3 --- Đặt $x=0$ và đổi biến $\tau=t-s$.} Tại $P=(0,0,0)$ cận trong là $\xi\in(-2(t-s),2(t-s))$. Đặt $\tau:=t-s$ (thời gian \textbf{đã trôi} kể từ khi sóng con phát đến lúc $t$). Khi $s:0\to t$ thì $\tau:t\to0$ và $ds=-d\tau$, nên
\[
  v_1(0,t)=\frac{1}{2c}\int_0^t\!\!\int_{-c\tau}^{c\tau}\!\chi_{[2024,2026]}(\xi)\,d\xi\,d\tau
  =\frac{1}{4}\int_0^t L(\tau)\,d\tau,
\]
với $L(\tau):=\big|[2024,2026]\cap[-2\tau,2\tau]\big|$ là \textbf{độ dài giao}.

\emph{Bước 1.4 --- Tính $L(\tau)$ (chia khúc).} Vì $[2024,2026]\subset(0,\infty)$, giao thực chất với $[0,2\tau]$:
\begin{itemize}
\item Nếu $2\tau<2024$ (tức $\tau<1012$): hai khoảng tách rời, $L(\tau)=0$.
\item Nếu $2024\le2\tau\le2026$ (tức $1012\le\tau\le1013$): giao $=[2024,2\tau]$, $L(\tau)=2\tau-2024$.
\item Nếu $2\tau>2026$ (tức $\tau>1013$): giao $=[2024,2026]$, $L(\tau)=2$.
\end{itemize}
\[
  L(\tau)=\begin{cases}0,&\tau<1012,\\ 2\tau-2024,&1012\le\tau\le1013,\\ 2,&\tau>1013.\end{cases}
\]

\emph{Bước 1.5 --- Tích phân theo $\tau$ (chia khúc đa thức).}
\begin{itemize}
\item $0\le t\le1012$: $L\equiv0$ trên $[0,t]\Rightarrow v_1(0,t)=0$.
\item $1012\le t\le1013$: $v_1(0,t)=\dfrac{1}{4}\displaystyle\int_{1012}^{t}(2\tau-2024)\,d\tau=\dfrac{1}{4}\big[(\tau-1012)^2\big]_{1012}^{t}=\dfrac{(t-1012)^2}{4}$.
\item $t\ge1013$: $v_1(0,t)=\dfrac{1}{4}\Big[\displaystyle\int_{1012}^{1013}(2\tau-2024)\,d\tau+\int_{1013}^{t}\!2\,d\tau\Big]=\dfrac{1+2(t-1013)}{4}=\dfrac{2t-2025}{4}$.
\end{itemize}
Kết luận:
\[
  u_1(0,0,0,t)=v_1(0,t)=\begin{cases}
    0,& 0\le t\le1012,\\[2pt]
    \dfrac{(t-1012)^2}{4},& 1012\le t\le1013,\\[6pt]
    \dfrac{2t-2025}{4},& t\ge1013.
  \end{cases}
\]

\paragraph{(b) Bước 2 --- $u_2$ (nguồn đối xứng cầu).}\leavevmode

\emph{Bước 2.1 --- Đặt biến bán kính.} Nguồn $f_2=\chi_{[0,1]}(t)\,\chi_B(x,y,z)$ \textbf{đối xứng cầu} quanh $O'=(2025,0,0)$ (vì $\chi_B$ chỉ phụ thuộc $\rho:=|(x,y,z)-O'|$). Vì dữ liệu đầu $=0$ cũng đối xứng cầu quanh $O'$, nghiệm $u_2$ cũng chỉ phụ thuộc $(\rho,t)$: viết $u_2(x,y,z,t)=U(\rho,t)$. Điểm hỏi $P=(0,0,0)$ cách $O'$ một khoảng $R=2025$, nên cuối cùng ta cần $U(R,t)$.

Trong toạ độ cầu, Laplacian của hàm chỉ phụ thuộc $\rho$ là $\Delta U=U_{\rho\rho}+\tfrac{2}{\rho}U_\rho$, nên $U$ thoả
\[
  U_{tt}=c^2\Big(U_{\rho\rho}+\tfrac{2}{\rho}U_\rho\Big)+\chi_{[0,1]}(t)\,\chi_{[0,1]}(\rho),\quad \rho>0,\ c=2,
\]
với $U(\rho,0)=U_t(\rho,0)=0$.

\emph{Bước 2.2 --- Đổi biến $v=\rho U$ để khử số hạng $\tfrac{2}{\rho}U_\rho$.} (Xem ``Đổi biến cầu $v=ru$''.) Đặt $v(\rho,t)=\rho\,U(\rho,t)$. Tính trực tiếp $v_{\rho\rho}=\rho\big(U_{\rho\rho}+\tfrac{2}{\rho}U_\rho\big)$, nên nhân phương trình trên với $\rho$ được
\[
  v_{tt}=c^2 v_{\rho\rho}+\rho\,\chi_{[0,1]}(t)\,\chi_{[0,1]}(\rho),\qquad \rho>0,
\]
với $v(\rho,0)=v_t(\rho,0)=0$. Đây là \textbf{sóng 1D có nguồn} trên nửa đường $\rho>0$.

\emph{Bước 2.3 --- Điều kiện biên tại tâm.} Vì $U$ phải hữu hạn ở tâm và $v=\rho U$, lấy $\rho\to0$ được
\[
  v(0,t)=0\qquad (\text{biên \textbf{Dirichlet} tự nhiên ở tâm}).
\]

\emph{Bước 2.4 --- Thác triển lẻ qua $\rho=0$ để dùng d'Alembert toàn trục.}
\begin{luuybox}
Bài toán đang trên \emph{nửa đường} $\rho>0$ với biên Dirichlet $v(0,t)=0$. Mẹo chuẩn: \textbf{thác triển lẻ} dữ liệu/nguồn sang $\rho<0$, rồi giải sóng 1D trên \emph{toàn trục} bằng d'Alembert. Vì hệ tuyến tính và dữ kiện đầu lẻ (ở đây $=0$, vừa chẵn vừa lẻ) cộng nguồn lẻ, nghiệm thu được cũng \emph{lẻ} theo $\rho$, do đó tự thoả $v(0,t)=0$. Quy tắc tổng quát:
\begin{itemize}
\item Biên Dirichlet ($v=0$) $\Rightarrow$ thác triển \textbf{lẻ} (hàm lẻ tự triệt tiêu tại $0$).
\item Biên Neumann ($v_\rho=0$) $\Rightarrow$ thác triển \textbf{chẵn} (hàm chẵn có đạo hàm $=0$ tại $0$).
\end{itemize}

Áp vào bài: nguồn trên $\rho>0$ là $g(\rho,s)=\rho\,\chi_{[0,1]}(s)\,\chi_{[0,1]}(\rho)$. Để $g$ trở thành hàm \emph{lẻ} theo $\rho$ trên cả $\mathbb{R}$, thác triển: $\rho$ vốn đã lẻ; còn $\chi_{[0,1]}(\rho)$ thác triển \emph{chẵn} thành $\chi_{[-1,1]}(\rho)$ (lẻ $\times$ chẵn = lẻ). Vậy nguồn mở rộng toàn trục là
\[
  \tilde g(\rho,s)=\rho\,\chi_{[0,1]}(s)\,\chi_{[-1,1]}(\rho).
\]
\end{luuybox}

\emph{Bước 2.5 --- Áp công thức Duhamel 1D cho $v$ trên toàn trục.} Với dữ kiện đầu $=0$ và nguồn $\tilde g$:
\[
  v(\rho,t)=\frac{1}{2c}\int_0^t\!\!\int_{\rho-c(t-s)}^{\rho+c(t-s)}\!\xi\,\chi_{[-1,1]}(\xi)\,\chi_{[0,1]}(s)\,d\xi\,ds.
\]
($\xi$ là biến tích phân chạy trên trục $\rho$; cận trong là miền phụ thuộc của $(\rho,t)$ qua một sóng con phát tại thời điểm $s$.)

\emph{Bước 2.6 --- Quay về $u_2$ tại $P$.} Vì $U=v/\rho$ và $P$ cách $O'$ một khoảng $R=2025$:
\[
  u_2(0,0,0,t)=U(R,t)=\frac{v(R,t)}{R}
  =\frac{1}{2cR}\iint_{\Delta(R,t)}\!\xi\,\chi_{[-1,1]}(\xi)\,\chi_{[0,1]}(s)\,d\xi\,ds,
\]
với $\Delta(R,t)=\{\,0<s<t,\ R-c(t-s)<\xi<R+c(t-s)\,\}$ và $\dfrac{1}{2cR}=\dfrac{1}{2\cdot2\cdot2025}=\dfrac{1}{8100}$.

\emph{Bước 2.7 --- Khoảng $t$ có tín hiệu \& tính tích phân.} Tín hiệu chỉ tới $P$ khi nón phụ thuộc giao được vùng nguồn $\{|\xi|\le1,\ s\in[0,1]\}$: cần $c(t-s)\ge R-1=2024$, kết hợp $s\in[0,1]$ cho $t-s\in[1012,1014]$, tức $t\in[1012,1014]$ (ngoài khoảng này $u_2=0$). Tính giao hình học (chia theo $t$) ra:
\[
  u_2(0,0,0,t)=\begin{cases}
    0,& 0\le t\le1012\ \text{hay}\ t\ge1014,\\[2pt]
    \dfrac{(t-1012)^2(2027-2t)}{3\cdot8100},& 1012\le t\le1013,\\[6pt]
    \dfrac{(t-1014)^2(2t-2025)}{3\cdot8100},& 1013\le t\le1014.
  \end{cases}
\]

\paragraph{(b) Bước 3 --- gộp $u=u_1-u_2$.}\leavevmode

\emph{Bước 3.1 --- Đồng bộ khoảng chia.} Cả $u_1$ và $u_2$ \emph{rời cân bằng} tại $t=1012$. Mốc còn lại:
\begin{itemize}
\item $u_1$ đổi khúc tại $t=1013$ (khi đó dải nguồn đã ``thấy'' đủ).
\item $u_2$ đổi khúc tại $t=1013$ và \emph{tắt} hẳn ($u_2=0$) khi $t\ge1014$ vì nguồn cầu chỉ phát trong $s\in[0,1]$.
\end{itemize}
Vậy chia $u=u_1-u_2$ theo bốn khúc: $[0,1012],\ [1012,1013],\ [1013,1014],\ [1014,\infty)$.

\emph{Bước 3.2 --- Trừ trên từng khúc.}
\begin{itemize}
\item $0\le t\le1012$: $u_1=u_2=0\Rightarrow u=0$.
\item $1012\le t\le1013$: $u_1=\dfrac{(t-1012)^2}{4}$, $u_2=\dfrac{(t-1012)^2(2027-2t)}{3\cdot8100}$. Đặt nhân tử chung $(t-1012)^2$:
\[
  u=(t-1012)^2\Big(\dfrac{1}{4}-\dfrac{2027-2t}{3\cdot8100}\Big).
\]
\item $1013\le t\le1014$: $u_1=\dfrac{2t-2025}{4}$, $u_2=\dfrac{(t-1014)^2(2t-2025)}{3\cdot8100}$. Đặt nhân tử chung $(2t-2025)$:
\[
  u=(2t-2025)\Big(\dfrac{1}{4}-\dfrac{(t-1014)^2}{3\cdot8100}\Big).
\]
\item $t\ge1014$: $u_2=0$, còn $u_1=\dfrac{2t-2025}{4}\Rightarrow u=\dfrac{2t-2025}{4}$.
\end{itemize}

\begin{nhobox}
\[
  u(0,0,0,t)=\begin{cases}
    0,& 0\le t\le1012,\\[2pt]
    (t-1012)^2\Big(\dfrac14-\dfrac{2027-2t}{3\cdot8100}\Big),& 1012\le t\le1013,\\[8pt]
    (2t-2025)\Big(\dfrac14-\dfrac{(t-1014)^2}{3\cdot8100}\Big),& 1013\le t\le1014,\\[8pt]
    \dfrac{2t-2025}{4},& t\ge1014.
  \end{cases}
\]
\textbf{Kiểm tra khớp:} (a) bằng $0$ tới $t=1012$ rồi rời cân bằng. Tại $t=1013$, hai biểu thức cho cùng giá trị $\dfrac{1}{4}-\dfrac{1}{3\cdot8100}$ (liên tục). Sau $t=1014$ quả cầu (nguồn tắt từ $s=1$) hết ảnh hưởng tới $P$, chỉ còn phần dải phẳng $u_1$.
\end{nhobox}

\subsection*{Dạng bài lân cận}
\begin{itemize}
\item \textbf{Nguồn phẳng đơn thuần} (hỏi tại điểm xa) --- xem ``đối xứng phẳng (hạ 1 biến)''.
\item \textbf{Nguồn/dữ liệu đối xứng cầu} --- xem ``đối xứng cầu ($v=ru$)''.
\item \textbf{Kirchhoff thuần nhất} (bỏ nguồn) --- xem ``Cauchy 3D thuần nhất''.
\end{itemize}
